\section{Randomized --- Uniform Sample}
%%% generated by the blueprint pipeline from TCSlib.Complexity.Randomized.UniformSample
\subsection*{Randomized --- Uniform Sample: Introduction}

Declarations of the Lean module \texttt{TCSlib.Complexity.Randomized.UniformSample}.

\subsection*{Randomized --- Uniform Sample: Results}

\begin{definition}[Bounded repeated sampling]
\lean{Randomized.sampleTries}
\label{Randomized.sampleTries}
\leanok
\uses{Complexity.LogProg.bitsVal}
For natural numbers $N,L$, a natural-number attempt budget $k$, and a Boolean list $r$,
define a partial result as follows. With no attempts, return no result; otherwise,
inspect the first $L$ bits of the current list, interpret them as a natural number with
the first bit least significant, and return that number if it is less than $N$. If it is
not, discard those $L$ bits and repeat with one fewer attempt; if no attempt succeeds,
return no result.
\end{definition}

\begin{definition}[Uniform sample procedure]
\lean{Randomized.uniformSample}
\statementsource{arora-barak-ch07-randomized-computation}{pdf-26d95080dbde-p003-b003}
\label{Randomized.uniformSample}
\leanok
\uses{Randomized.sampleTries}
For natural numbers $N$ and $k$ and a Boolean list $r$, the output is the result of
applying the sample-tries procedure with parameters $N$, $\operatorname{size}(N-1)$,
$k$, and $r$.
\end{definition}

\begin{definition}[Number of sample bits]
\lean{Randomized.sampleCoins}
\statementsource{arora-barak-ch07-randomized-computation}{
  pdf-26d95080dbde-p002-b002,
  pdf-26d95080dbde-p003-b003
}
\label{Randomized.sampleCoins}
\leanok
For natural numbers $N$ and $k$, the number of sample bits is
$k\operatorname{size}(N-1)$.
\end{definition}

\begin{theorem}[Repeated sampling as a list search]
\lean{Randomized.sampleTries_eq_find}
\label{Randomized.sampleTries_eq_find}
\leanok
\uses{Complexity.LogProg.bitsVal, Randomized.blockVal, Randomized.sampleTries}
\difficulty{4}
Let $N,L,k$ be natural numbers and let $r$ be a finite list of bits. The result of
sampling at most $k$ successive length-$L$ blocks of $r$, interpreting each block as a
binary number with its first bit least significant, is the first value less than $N$ in
the resulting list, or no value if there is none.
\end{theorem}

\begin{theorem}[Range of a uniform sample]
\lean{Randomized.lt_of_uniformSample_eq_some}
\label{Randomized.lt_of_uniformSample_eq_some}
\leanok
\uses{Randomized.sampleTries_eq_find, Randomized.uniformSample}
\difficulty{2}
For natural numbers $N,k$, a bit string $r$, and a natural number $i$, if the sampling
procedure with parameters $N,k$ and input $r$ returns $i$, then $i<N$.
\end{theorem}

\begin{theorem}[Bounds from binary size]
\lean{Randomized.le_two_pow_size_pred}
\label{Randomized.le_two_pow_size_pred}
\leanok
\difficulty{4}
For every natural number $N$ with $1\leq N$, the power of two
$2^{\operatorname{size}(N-1)}$ is at least $N$ and is strictly less than $2N$.
\end{theorem}

\begin{theorem}[Probability of Repeated Uniform Sampling]
\lean{Randomized.randProb_uniformSample}
\label{Randomized.randProb_uniformSample}
\leanok
\uses{Randomized.randProb, Randomized.randProb_blocks_piFinset, Randomized.randProb_congr, Randomized.sampleCoins, Randomized.sampleTries_eq_find, Randomized.uniformSample}
\difficulty{3}
Let $N,k\in\mathbb{N}$, and let $P$ be any property of optional natural numbers. Put
$L=\mathrm{size}(N-1)$, where subtraction is truncated at zero. The probability that the
procedure which reads $k$ consecutive blocks of $L$ random bits, returns the first block
value less than $N$, and returns no value if there is none satisfying this condition,
produces an output satisfying $P$ equals $1/2^{kL}$ times the number of tuples
$(a_1,\ldots,a_k)\in\{0,\ldots,2^L-1\}^k$ for which the first $a_t<N$ exists and $P$
holds of that first such value, or for which no such $a_t$ exists and $P$ holds of no
value.
\end{theorem}

\begin{theorem}[Equal Probabilities for Accepted Outputs]
\lean{Randomized.randProb_uniformSample_eq}
\label{Randomized.randProb_uniformSample_eq}
\leanok
\uses{Randomized.le_two_pow_size_pred, Randomized.randProb, Randomized.randProb_uniformSample, Randomized.sampleCoins, Randomized.uniformSample}
\difficulty{5}
Let $N,k,i,j\in\mathbb{N}$ with $i<N$ and $j<N$. Under the repeated sampling procedure
that reads $k$ blocks of random bits and returns the first block value less than $N$,
the probability of returning $i$ equals the probability of returning $j$.
\end{theorem}

\begin{theorem}[Probability of No Accepted Sample]
\lean{Randomized.randProb_uniformSample_none}
\label{Randomized.randProb_uniformSample_none}
\leanok
\uses{Randomized.blockVal, Randomized.randProb, Randomized.randProb_bitsVal, Randomized.randProb_congr, Randomized.randProb_forall_blocks, Randomized.sampleCoins, Randomized.sampleTries_eq_find, Randomized.uniformSample}
\difficulty{4}
Let $N,k\in\mathbb{N}$ and put $L=\mathrm{size}(N-1)$, with subtraction truncated at
zero. In the procedure that reads $k$ blocks of $L$ random bits and accepts a block when
its value is less than $N$, the probability that no block is accepted is
$\left(\frac{2^L-N}{2^L}\right)^k$.
\end{theorem}

\begin{theorem}[Bound on Sampling Failure]
\lean{Randomized.randProb_uniformSample_none_le}
\label{Randomized.randProb_uniformSample_none_le}
\leanok
\uses{Randomized.le_two_pow_size_pred, Randomized.randProb, Randomized.randProb_uniformSample_none, Randomized.sampleCoins, Randomized.uniformSample}
\difficulty{4}
Let $N,k\in\mathbb{N}$ with $1\le N$. In the procedure that reads $k$ blocks of random
bits and returns the first block value less than $N$, the probability that no block is
accepted is at most $(1/2)^k$.
\end{theorem}

\begin{theorem}[Accepted Output Probability Relation]
\lean{Randomized.randProb_uniformSample_some_mul}
\label{Randomized.randProb_uniformSample_some_mul}
\leanok
\uses{Randomized.lt_of_uniformSample_eq_some, Randomized.randProb, Randomized.randProb_congr, Randomized.randProb_mem_eq_sum, Randomized.randProb_true, Randomized.randProb_uniformSample_eq, Randomized.sampleCoins, Randomized.uniformSample}
\difficulty{5}
Let $N,k,i\in\mathbb{N}$ with $i<N$. In the repeated sampling procedure, the probability
of returning $i$, multiplied by $N$, equals one minus the probability of returning no
value.
\end{theorem}

\begin{definition}[Encoding an optional natural number]
\lean{Randomized.encSample}
\label{Randomized.encSample}
\leanok
For an optional natural number, define its Boolean-list encoding by sending `none` to
the empty list and `some v` to the list whose first entry is true and whose remaining
entries are the binary digits of $v$.
\end{definition}

\begin{definition}[Sampling-state update]
\lean{Randomized.sampleStep}
\label{Randomized.sampleStep}
\leanok
\uses{Complexity.BinArith.lt, Complexity.LogProg.bitsVal, Complexity.isNilB, Complexity.pairFstD, Complexity.pairSndD, Randomized.dropPrefixByLen, Randomized.takePrefixByLen, Turing.pairEncode}
Define a transformation of a Boolean list by decoding its nested paired data into a
first component, a candidate bit block, and a remaining block. The output pairs the
first component with the unused suffix of the candidate block; if the remaining block is
empty and the candidate block represents a number smaller than the encoded bound, append
a true bit and that number's binary digits, and otherwise retain the remaining block.
\end{definition}

\begin{definition}[Uniform sampling algorithm]
\lean{Randomized.uniformSampleAlg}
\label{Randomized.uniformSampleAlg}
\leanok
\uses{Complexity.BinArith.forLoop, Complexity.BinArith.sub, Complexity.pairFstD, Complexity.pairSndD, Randomized.sampleStep, Turing.pairEncode}
The uniform sampling algorithm is the function on Boolean lists obtained by applying the
specified loop with step function `sampleStep` to the encoded components of its input
and returning the second component of the resulting encoding.
\end{definition}

\begin{theorem}[Iteration of the sampling step]
\lean{Randomized.iterate_sampleStep}
\label{Randomized.iterate_sampleStep}
\leanok
\uses{Complexity.BinArith.lt, Complexity.LogProg.bitsVal, Complexity.isNilB, Complexity.pairFstD_pairEncode, Complexity.pairSndD_pairEncode, Randomized.dropPrefixByLen, Randomized.encSample, Randomized.sampleStep, Randomized.sampleTries, Randomized.takePrefixByLen, Turing.pairEncode}
\difficulty{4}
For all bit strings $Nw$, $Lw$, $r$, and $out$, and every natural number $j$, iterating
the sampling step $j$ times on the encoding of the pair $(Nw,Lw)$ and the pair $(r,out)$
gives the encoding of $(Nw,Lw)$ and a pair whose first component is obtained from $r$ by
deleting its first $j\,|Lw|$ bits. Its second component is $out$ if $out$ is nonempty;
if $out$ is empty, it is the encoding of the sample returned by making at most $j$
sampling tries with threshold given by the binary value of $Nw$, using blocks of length
$|Lw|$ from $r$.
\end{theorem}

\begin{theorem}[Correctness of the uniform sampling algorithm]
\lean{Randomized.uniformSampleAlg_pairEncode}
\label{Randomized.uniformSampleAlg_pairEncode}
\leanok
\uses{Complexity.BinArith.forLoop_pairEncode, Complexity.BinArith.sub, Complexity.LogProg.bitsVal, Complexity.pairFstD_pairEncode, Complexity.pairSndD_pairEncode, Randomized.encSample, Randomized.iterate_sampleStep, Randomized.uniformSample, Randomized.uniformSampleAlg, Turing.pairEncode}
\difficulty{3}
For bit strings $w,r$ and a natural number $k$, encode the pair consisting of the
encoding of $w$ with $k$ copies of the bit $1$ and the string $r$. The uniform-sampling
algorithm on this input returns the empty string if none of the first $k$ consecutive
blocks of length $\operatorname{size}(N-1)$ in $r$ has binary value less than $N$, where
$N$ is the binary value of $w$; otherwise it returns a leading bit $1$ followed by the
binary expansion of the first such block's value. In the displayed length,
$\operatorname{size}$ is the binary length used by the algorithm.
\end{theorem}

\begin{theorem}[Polynomial-time sampling step]
\lean{Randomized.polyTimeComputable_sampleStep}
\label{Randomized.polyTimeComputable_sampleStep}
\leanok
\uses{Complexity.BinArith.lt, Complexity.BinArith.polyTimeComputable_lt, Complexity.BinArith.polyTimeComputable_normBits, Complexity.BinArith.polyTimeComputable_slot2, Complexity.BinArith.polyTimeComputable_slot3, Complexity.LogProg.bitsVal, Complexity.PolyTimeComputable, Complexity.PolyTimeComputable.append, Complexity.PolyTimeComputable.comp, Complexity.PolyTimeComputable.pairEncode, Complexity.isNilB, Complexity.pairFstD, Complexity.pairSndD, Complexity.polyTimeComputable_and, Complexity.polyTimeComputable_const, Complexity.polyTimeComputable_isNil, Complexity.polyTimeComputable_ite, Complexity.polyTimeComputable_pairFstD, Complexity.polyTimeComputable_pairSndD, Randomized.polyTimeComputable_dropPrefixByLen, Randomized.polyTimeComputable_takePrefixByLen, Randomized.sampleStep, Randomized.takePrefixByLen, Turing.pairEncode}
\difficulty{4}
The bit-string function that takes a state encoding, splits its encoded sample into a
prefix of length given by the encoded state parameter and the remaining suffix, and
updates the state and sample according to the specified comparison and binary-value
rules is polynomial-time computable.
\end{theorem}

\begin{theorem}[Polynomial-time uniform-sampling algorithm]
\lean{Randomized.polyTimeComputable_uniformSampleAlg}
\label{Randomized.polyTimeComputable_uniformSampleAlg}
\leanok
\uses{Complexity.BinArith.forLoop_bound_of_ctx, Complexity.BinArith.length_bits_bitsVal_le, Complexity.BinArith.polyTimeComputable_forLoop, Complexity.BinArith.polyTimeComputable_slot4, Complexity.BinArith.polyTimeComputable_sub, Complexity.BinArith.sub, Complexity.PolyTimeComputable, Complexity.PolyTimeComputable.comp, Complexity.PolyTimeComputable.pairEncode, Complexity.length_pairSndD_le, Complexity.length_pair_components_le, Complexity.pairFstD, Complexity.pairFstD_pairEncode, Complexity.pairSndD, Complexity.pairSndD_pairEncode, Complexity.polyTimeComputable_const, Complexity.polyTimeComputable_pairFstD, Complexity.polyTimeComputable_pairSndD, Randomized.dropPrefixByLen, Randomized.polyTimeComputable_sampleStep, Randomized.sampleStep, Randomized.takePrefixByLen, Randomized.uniformSampleAlg, Turing.length_pairEncode, Turing.pairEncode}
\difficulty{5}
The function that initializes the encoded sampling state from the decoded input and
iterates the sampling step for the encoded number of rounds is polynomial-time
computable; its output is the final sample component of the resulting state.
\end{theorem}

\begin{theorem}[Uniform sampling by repeated rejection]
\lean{Randomized.uniformSample_spec}
\label{Randomized.uniformSample_spec}
\leanok
\uses{Complexity.LogProg.bitsVal_bits, Complexity.PolyTimeComputable, Randomized.encSample, Randomized.lt_of_uniformSample_eq_some, Randomized.polyTimeComputable_uniformSampleAlg, Randomized.randProb, Randomized.randProb_uniformSample_none_le, Randomized.randProb_uniformSample_some_mul, Randomized.sampleCoins, Randomized.uniformSample, Randomized.uniformSampleAlg, Randomized.uniformSampleAlg_pairEncode, Turing.pairEncode}
\difficulty{4}
Let $N\ge 1$ and let $\delta>0$. Set $k=\lceil\log_2\lceil\delta^{-1}\rceil\rceil$ and
$m=k\,\mathrm{size}(N-1)$. The algorithm defining uniform sampling is computable in
polynomial time; on input the pair encoding of the pair consisting of the binary
representation of $N$ followed by $k$ ones, and an arbitrary bit string $r$, it returns
the encoding of the result of sampling from $r$ with parameters $N,k$. Every returned
value is less than $N$. Under the uniform distribution on bit strings of length $m$, for
each $i<N$, the probability of returning $i$, multiplied by $N$, equals one minus the
probability of returning no value; moreover, the probability of returning no value is at
most $\delta$.
\end{theorem}
